% =============================================================================
%  problem_fit.tex -- Section: Problem Fit, User Pain Points, and AI Rationale
%  BodyCheck -- AI-Powered Medical Triage & 3D Symptom Checker
%  Team Clavicular | LotusHacks 2026
%
%  Pure content module -- include via \input{problem_fit}.
%  Uses only packages already loaded in main.tex.
% =============================================================================

\section{Problem Fit, User Pain Points, and AI Rationale}
\label{sec:problem-fit}

This section addresses the rubric criterion on \textit{tinh phu hop de bai}
--- alignment with the problem statement, evidence-based analysis of end-user
pain points, and a clear explanation of why artificial intelligence is the
appropriate technical approach. BodyCheck directly targets early medical
triage: users localise pain on an interactive body model, provide symptoms
through text or voice, receive uncertainty-aware first-pass guidance, and are
directed toward appropriate next steps or nearby healthcare facilities.

The core problem is not simply that users lack medical information. The more
specific problem is that non-specialist users often face a gap between symptom
onset and professional care. During that gap, they must decide whether a symptom
is manageable at home, requires a routine appointment, or may require urgent
care. Current alternatives --- search engines, social media, static symptom
checklists, and manual triage calls --- are either unreliable, too slow, too
expensive, or too dependent on clinical expertise. BodyCheck is therefore
designed as a triage-support tool, not as a final diagnostic authority.

% ---------------------------------------------------------------------------
% Table 1 -- User Pain-Point Evidence
% ---------------------------------------------------------------------------
\begin{table}[htbp]
  \centering
  \caption{End-user pain points, supporting evidence, and BodyCheck response.}
  \label{tab:problem-fit-pain-points}
  \setlength{\tabcolsep}{5pt}
  \renewcommand{\arraystretch}{1.45}
  \small
  \begin{tabular}{@{} p{2.9cm} p{4.2cm} p{4.2cm} p{2.2cm} @{}}
    \toprule
    \textbf{Pain Point} &
    \textbf{Observed Evidence} &
    \textbf{BodyCheck Response} &
    \textbf{Fit} \\
    \midrule

    Uneven access to timely care &
    Vietnam remains a mixed urban--rural healthcare market: DataReportal
    reports that 59.5\% of Vietnam's population lived in rural areas in
    early 2025.\footnotemark[1] &
    Provides immediate first-pass triage through a web/mobile interface,
    reducing the need for a user to wait until clinic access is available
    before understanding the likely urgency of their symptoms. &
    \textcolor{clav-green}{\textbf{High}} \\

    Slow early triage &
    Vietnam had 12.5 doctors per 10{,}000 people in 2023, according to a
    Ministry-of-Health-reported figure cited by VietnamNet.\footnotemark[2]
    This indicates that clinician time remains a scarce resource relative
    to population-scale triage demand. &
    Automates the first structured intake step: body region, symptom text,
    follow-up questions, risk level, and suggested next action. Clinicians
    can receive a cleaner summary instead of raw, unstructured complaints. &
    \textcolor{clav-green}{\textbf{High}} \\

    Risky self-diagnosis behaviour &
    Pew Research Center found that 35\% of U.S. adults had gone online to
    identify a possible medical condition, while 18\% of online diagnosers
    later received a different or non-matching opinion from a clinician.
    \footnotemark[3] &
    Replaces ad-hoc searching with a constrained triage workflow:
    uncertainty wording, safety disclaimers, red-flag escalation, and
    references rather than unsupported diagnostic certainty. &
    \textcolor{clav-green}{\textbf{High}} \\

    Low-friction digital expectation &
    Vietnam had 79.8 million internet users in January 2025, equal to
    78.8\% internet penetration.\footnotemark[1] This supports a
    mobile/web-first access model for consumer health guidance. &
    Uses Flutter/web delivery, body-map interaction, text or voice symptom
    input, and short response formats so non-specialist users do not need
    medical vocabulary to begin the triage flow. &
    \textcolor{clav-green}{\textbf{High}} \\

    Need for understandable next steps &
    The pitch-deck user flow requires users to move from body-map selection
    to symptoms, AI reasoning, possible conditions, recommended actions,
    references, and nearby clinics. &
    BodyCheck maps the full decision path: ``What hurts?'', ``What symptoms
    accompany it?'', ``How urgent is this?'', ``What should I do next?'',
    and ``Where can I seek care?'' &
    \textcolor{clav-green}{\textbf{High}} \\

    \bottomrule
  \end{tabular}

  \smallskip\noindent
  {\footnotesize
  \footnotemark[1]\href{https://datareportal.com/reports/digital-2025-vietnam}
  {DataReportal, \textit{Digital 2025: Vietnam}.}
  \quad
  \footnotemark[2]\href{https://vietnamnet.vn/en/vn-targets-15-doctors-per-10-000-people-by-2025-2254601.html}
  {VietnamNet, citing Ministry of Health figures, 2024.}
  \quad
  \footnotemark[3]\href{https://www.pewresearch.org/internet/2013/01/15/health-online-2013/}
  {Pew Research Center, \textit{Health Online 2013}.}
  }
\end{table}

% =============================================================================
\subsection{Alignment with the Project Requirement}
\label{sec:problem-fit-alignment}
% =============================================================================

BodyCheck is aligned with the problem requirement because it focuses on a
concrete, high-frequency user decision: how to respond when pain or discomfort
appears but professional medical help is not immediately available. The product
does not attempt to replace a physician. Instead, it supports the earlier and
safer step of triage: collecting symptoms, identifying urgency signals, giving
plain-language guidance, and encouraging escalation when needed.

The attached project outline defines the product around three linked actions:
localising pain through a 3D body model, generating clinician-style preliminary
guidance, and connecting users with relevant healthcare facilities. This is a
direct match for users who cannot easily convert a vague sensation such as
``chest tightness'', ``shoulder pain'', or ``abdominal discomfort'' into a clear
care pathway. The body-map interaction reduces ambiguity at input time, while
the AI layer transforms the user's natural-language description into a
structured response.

The pitch flow further reinforces the fit: users tap the body map, add symptoms
through text or voice, receive AI reasoning and risk scoring, read conditions
and recommended actions, and then find nearby treatment options. This is not an
abstract AI feature layered onto a healthcare product; it is the central
mechanism that converts unstructured symptoms into actionable triage guidance.

% =============================================================================
\subsection{Why AI Is Necessary}
\label{sec:problem-fit-why-ai}
% =============================================================================

AI is necessary because the input problem is inherently unstructured. Users do
not describe symptoms in standard medical terminology. They may mix location,
duration, severity, lifestyle context, medication history, and emotional concern
in one short message. A static rule tree can handle only a limited set of
predefined cases, while a search engine returns documents that the user still
has to interpret. BodyCheck requires an intermediate reasoning layer that can
parse free text, ask follow-up questions, communicate uncertainty, and produce
a structured output that is understandable to both users and downstream systems.

\begin{table}[htbp]
  \centering
  \caption{Why the BodyCheck use case requires AI rather than a purely static workflow.}
  \label{tab:why-ai}
  \setlength{\tabcolsep}{5pt}
  \renewcommand{\arraystretch}{1.45}
  \small
  \begin{tabular}{@{} p{3.0cm} p{4.0cm} p{4.4cm} p{2.0cm} @{}}
    \toprule
    \textbf{Required Capability} &
    \textbf{Non-AI Limitation} &
    \textbf{AI Contribution} &
    \textbf{Safety Control} \\
    \midrule

    Natural symptom understanding &
    Users write vague or incomplete descriptions that do not match exact
    checklist labels. &
    Extracts body region, symptom type, duration, severity, and contextual
    clues from free text. &
    Structured prompt template and validated response schema. \\

    Follow-up questioning &
    Static forms either ask too many questions or miss case-specific details. &
    Generates targeted follow-up questions based on the selected body region
    and symptoms already provided. &
    Maximum question count and red-flag escalation rules. \\

    Urgency guidance &
    Search results do not reliably distinguish home care, routine care, and
    urgent care. &
    Produces a triage-style risk level with recommended next steps and
    uncertainty-aware explanation. &
    Mandatory disclaimer and emergency escalation language. \\

    Evidence summarisation &
    Users cannot quickly compare long medical references during a stressful
    moment. &
    Retrieves, ranks, and summarises relevant sources in plain language. &
    Source links are preserved; definitive diagnosis is prohibited. \\

    Multilingual and accessible output &
    Text-heavy interfaces exclude users who prefer voice or simpler language. &
    Supports voice input/output and can adapt response phrasing for English
    and Vietnamese users. &
    Medical meaning is constrained by response format and clinical disclaimer. \\

    \bottomrule
  \end{tabular}
\end{table}

The strongest reason for using AI is therefore not novelty. The strongest reason
is task fit. BodyCheck must mediate between messy human symptom descriptions and
structured triage outputs. AI is appropriate because it can handle linguistic
variation, maintain conversational context, summarise evidence, and produce
consistent first-pass guidance. At the same time, the system deliberately
limits the role of AI: it does not issue final diagnoses, does not replace
clinical judgement, and must escalate high-risk symptoms to professional or
emergency care.

% =============================================================================
\subsection{End-User Value Proposition}
\label{sec:problem-fit-value}
% =============================================================================

For health-conscious consumers, BodyCheck reduces confusion during the first
minutes after symptom onset. Instead of searching many unrelated pages, the user
can identify a body region, describe symptoms naturally, and receive a concise
summary of possible explanations and next steps. This is especially useful for
users who do not know medical vocabulary or who need guidance outside clinic
hours.

For clinical triage operators and telehealth nurses, BodyCheck reduces the
amount of unstructured intake work. The system can prepare a cleaner symptom
summary, identify missing information, and highlight possible red flags before
human review. The value is not replacing the operator; the value is making the
first-pass intake faster and more consistent.

For primary-care clinicians, BodyCheck can act as a pre-consultation summary
tool. A patient-facing intake flow that captures body region, symptoms,
duration, severity, and AI-generated references can reduce the time spent
clarifying basic symptom history during the appointment.

% ---------------------------------------------------------------------------
% Table 3 -- Requirement Coverage Matrix
% ---------------------------------------------------------------------------
\begin{table}[htbp]
  \centering
  \caption{Coverage of the 25-point ``problem fit'' requirement.}
  \label{tab:problem-fit-coverage}
  \setlength{\tabcolsep}{6pt}
  \renewcommand{\arraystretch}{1.4}
  \small
  \begin{tabular}{@{} p{4.0cm} p{7.2cm} p{2.0cm} @{}}
    \toprule
    \textbf{Rubric Requirement} &
    \textbf{How This Section Addresses It} &
    \textbf{Status} \\
    \midrule

    Closely follows the problem statement &
    Focuses on early medical triage, pain localisation, symptom intake,
    urgency guidance, references, and clinic handoff. &
    \textcolor{clav-green}{\textbf{Covered}} \\

    Analyses customer/end-user pain points &
    Identifies uneven access, slow triage, risky self-diagnosis,
    low-friction digital expectations, and need for understandable next steps. &
    \textcolor{clav-green}{\textbf{Covered}} \\

    Uses real data &
    Uses external public statistics on Vietnam digital readiness,
    rural population share, doctor density, and online self-diagnosis behaviour. &
    \textcolor{clav-amber}{\textbf{Partial}} \\

    Clearly explains ``why AI'' &
    Shows that free-text symptom understanding, follow-up questioning,
    urgency classification, evidence summarisation, and voice accessibility
    require AI-supported reasoning rather than a static checklist alone. &
    \textcolor{clav-green}{\textbf{Covered}} \\

    Maintains medical safety framing &
    States that BodyCheck supports triage only, avoids final diagnosis,
    communicates uncertainty, and escalates red-flag cases. &
    \textcolor{clav-green}{\textbf{Covered}} \\

    \bottomrule
  \end{tabular}

  \smallskip\noindent
  {\footnotesize
  The ``real data'' criterion is marked as partial because the current
  project files do not include primary user-research data from the team.
  To make this criterion fully defensible, the team should add its own
  survey, interview, or prototype-testing results before final submission.
  }
\end{table}

% =============================================================================
\subsection{Additional Primary Data Needed Before Submission}
\label{sec:problem-fit-data-needed}
% =============================================================================

The current section is strong enough as a public-data-backed argument, but the
rubric specifically rewards analysis of customer or end-user pain points. For
maximum score, the team should add at least one primary-data block from its own
validation work. The most useful additions are:

\begin{itemize}[leftmargin=2em, itemsep=3pt, topsep=4pt]
  \item \textbf{User survey:} number of respondents, percentage who have
        searched symptoms online, percentage who felt uncertain about urgency,
        and percentage who would use a body-map triage assistant.
  \item \textbf{Prototype test:} number of testers, average time to complete
        symptom input, task-completion rate, and qualitative comments on the
        body-map interaction.
  \item \textbf{Interview evidence:} three to five short pain-point findings
        from students, parents, rural users, telehealth users, or healthcare
        support staff.
  \item \textbf{Before--after comparison:} time required to search symptoms
        manually versus time required to complete the BodyCheck flow.
\end{itemize}

Once those numbers are available, Table~\ref{tab:problem-fit-pain-points} should
be updated so at least one row cites team-collected data rather than only
public secondary sources.

% =============================================================================
% End of Section: Problem Fit, User Pain Points, and AI Rationale
% =============================================================================